\documentclass{amsbook}

\usepackage{amssymb}
\usepackage{hyperref}
\usepackage{enumitem}
\usepackage{mathtools}
\usepackage{geometry}                % See geometry.pdf to learn the layout options. There are lots.
\geometry{a4paper}                   % ... or a4paper or a5paper or ... 

\newtheorem{theorem}{Theorem}[chapter]
\newtheorem{lemma}[theorem]{Lemma}

\theoremstyle{definition}
\newtheorem{definition}[theorem]{Definition}
\newtheorem{example}[theorem]{Example}
\newtheorem{xca}[theorem]{Exercise}

\theoremstyle{remark}
\newtheorem{remark}[theorem]{Remark}

\numberwithin{section}{chapter}
\numberwithin{equation}{chapter}

\begin{document}

\frontmatter

\title{Minqi's Solutions to the Exercises of Rick Durrett (2019), Probability: Theory and Examples (5th Edition)}

\author{Minqi Pan\footnote{Website: \href{http://www.minqi-pan.com}{http://www.minqi-pan.com}}}

\date{\today}

\maketitle

\setcounter{page}{4}

\tableofcontents

\mainmatter

\chapter{Measure Theory}

\section{Exercise 1.1.1: An Example of Probability Space}
Let $\Omega=\mathbb{R},\mathcal{F}=$ all subsets so that $A$ or $A^c$ is countable, $P(A)=0$ in the first case and $=1$ in the second. Show that $(\Omega,\mathcal{F},P)$ is a probability space.
\begin{proof}
First, we prove that $\mathcal F$ is a $\sigma-$algebra on $\mathbb{R}$.
\begin{enumerate}
\item $\mathcal F$ is nonempty because, say, $\{\sqrt{2}\}$ is in $\mathcal F$.
\item Suppose $A\in\mathcal F$, then, by definition of $\mathcal F$, $A=(A^c)^c$ or $A^c$ is countable. Hence $A^c\in\mathcal F$.
\item Suppose $A_1,A_2,\dots\in\mathcal{F}$ is a countable sequence of sets. Define\[
\begin{split}
\mathcal{A}&\equiv\{A_1,A_2,\dots\},\\
\mathcal{B}&\equiv\{X\in\mathcal{A}:X\text{ is countable}\},\\
\mathcal{C}&\equiv\mathcal{A}-\mathcal{B}.
\end{split}
\]If $\mathcal{C}=\varnothing$, then $\cup_{i=1}^\infty A_i$ is countable because of Theorem~2.12 of \cite{rudin1976principles}. Otherwise, pick any $A_m\in\mathcal{C}$, then $A_m^c$ must be countable, because $A_m\in\mathcal{F}$ and $A_m\notin\mathcal{B}$. Since\[
\left(\bigcup_{i=1}^\infty A_i\right)^c=\bigcap_{i=1}^\infty A_i^c\subset A_m^c,
\]Theorem~2.8 of \cite{rudin1976principles} implies that $(\cup_{i=1}^\infty A_i)^c$ is countable. Therefore $\cup_{i=1}^\infty A_i\in\mathcal F$.
\end{enumerate}

Next, we prove that $P$ is a measure on $(\mathbb{R},\mathcal{F})$.
\begin{enumerate}
\item $P$ is nonnegative because $P$ maps to either $0$ or $1$. $P(\varnothing)=0$ because $\varnothing$ is countable.
\item Suppose $A_1,A_2,\dots\in\mathcal{F}$ is a countable sequence of disjoint sets. Define $\mathcal{A}, \mathcal{B}, \mathcal{C}$ as above. If $\mathcal{C}=\varnothing$, then, as we have shown above, $\cup_{i=1}^\infty A_i$ is countable. Hence\[
P\left(\bigcup_{i=1}^\infty A_i\right)=0=\sum_{A_i\in\mathcal{B}} P(A_i)=\sum_{i=1}^\infty P(A_i).
\]Otherwise, we claim that $\mathcal{C}$ cannot have two distinct elements. To show this, suppose $C_1,C_2\in\mathcal{C}$ and $C_1\ne C_2$. Then $C_1^c\cup C_2^c=\mathbb{R}$ because $C_1\cap C_2=\varnothing$. Then Theorem~2.8 of \cite{rudin1976principles} implies that $\mathbb{R}$ is countable, but this contradicts with the corollary of Theorem~2.43 of \cite{rudin1976principles}.Therefore $\mathcal{C}$ contains only one element, say, $A_m$. Hence\[
P\left(\bigcup_{i=1}^\infty A_i\right)=1=P(A_m)=\sum_{A_i\in\mathcal{B}} P(A_i)+\sum_{A_i\in\mathcal{C}} P(A_i)=\sum_{i=1}^\infty P(A_i).
\]
\end{enumerate}

Finally, we prove that $P$ is a probability measure. Note that $\mathbb{R}^c=\varnothing$ is countable. Hence\[
P(\mathbb{R})=1.
\]
\end{proof}
\section{Exercise 1.1.4: The Union of Increasing $\sigma-$algebras is NOT a $\sigma-$algebra}
\begin{enumerate}
\item Show that if $\mathcal{F}_1\subset\mathcal{F}_2\subset\dots$ are $\sigma-$algebras, then $\cup_i\mathcal{F}_i$ is an algebra.
\item Give an example to show that $\cup_i\mathcal{F}_i$ need not be a $\sigma-$algebra.
\end{enumerate}
\begin{proof}
\begin{enumerate}
\item Pick any $F\in\cup_i\mathcal{F}_i$, then there exists $i_0$ such that $F\in\mathcal{F}_{i_0}$. Since $\mathcal{F}_{i_0}$ is an algebra, $F^c\in\mathcal{F}_{i_0}$. Hence $F^c\in\cup_i\mathcal{F}_i$.

Further more, pick any $E,F\in\cup_i\mathcal{F}_i$, then there exists $i_0$ and $j_0$ such that $F\in\mathcal{F}_{i_0}$ and $E\in\mathcal{F}_{j_0}$. Without loss of generality, suppose $i_0\leqslant j_0$. Then $F\in\mathcal{F}_{j_0}$ because $\mathcal{F}_{i_0}\subset\mathcal{F}_{j_0}$. Since $F\in\mathcal{F}_{j_0}$ is an algebra, $E\cup F\in\mathcal{F}_{j_0}$. Therefore $E\cup F\in\cup_i\mathcal{F}_i$.

We have shown above that $\cup_i\mathcal{F}_i$ is closed under complements and pairwise unions, therefore it is an algebra.
\item Let $\Omega\equiv[0,1]$. Define\[
\begin{split}
\mathcal{F}_1&\equiv\sigma\left([0.0,0.1],[0.1,1.0]\right),\\
\mathcal{F}_2&\equiv\sigma\left([0.00,0.01],[0.01,0.10],[0.10,0.11],[0.11,1.00]\right),\\
&\dots.
\end{split}
\]where $0.1$ is the binary representation of $\frac{1}{2}$, $0.01$ is the binary representation of $\frac{1}{4}$, etc. Then it is clear that\[
\mathcal{F}_1\subset\mathcal{F}_2\subset\dots.
\]Pick an irrational number $x\in(0,1)$, and suppose its binary representation is\[
0.s_1s_2s_3\dots.
\]Then\[
[x,1]=\bigcap_{i=1}^\infty[0.s_1s_2\dots s_{i-1}s_i,1].
\]To arrive at a contradiction, suppose $\cup_i\mathcal{F}_i$ is a $\sigma-$algebra, then\[
[x,1]\in\cup_i\mathcal{F}_i,
\]because $[0.s_1s_2\dots s_{i-1}s_i,1]\in\cup_i\mathcal{F}_i$ for all $i$. Then there exists $i_0$ such that\[
[x,1]\in\mathcal{F}_{i_0},
\]which contradicts with the obvious fact that $\mathcal{F}_{i_0}$ only contains intervals of rational endpoints.
\end{enumerate}
\end{proof}

\section{Exercise 1.2.1: New Random Variable from Two Existing Random Variables}
Suppose $X$ and $Y$ are random variables on $(\Omega,F,P)$ and let $A\in\mathcal{F}$. Show that if we let $Z(\omega)=X(\omega)$ for $\omega\in A$ and $Z(\omega)=Y(\omega)$ for $\omega\in A^c$, then $Z$ is a random variable.
\begin{proof}
Pick any $\omega\in\Omega$, then $\omega\in A$ or $\omega\in A^c$. Hence $Z$ is defined on all $\omega\in\Omega$. Since both of $X$ and $Y$ map into $\mathbb{R}$, $Z$ maps to $\mathbb{R}$. Therefore $Z$ is a well-defined function from $\Omega$ to $\mathbb{R}$.

Next, we show that $Z$ is a measurable map. Pick any $B\in\mathcal{R}$, then\[
\begin{split}
Z^{-1}(B)&=\{\omega\in\Omega:Z(\omega)\in B\}\\
&=\{\omega\in A:Z(\omega)\in B\}\cup\{\omega\in A^c:Z(\omega)\in B\}\\
&=\{\omega\in A:X(\omega)\in B\}\cup\{\omega\in A^c:Y(\omega)\in B\}\\
&=(A\cap\{\omega\in\Omega:X(\omega)\in B\})\cup(A^c\cap\{\omega\in\Omega:Y(\omega)\in B\})\\
&=(A\cap X^{-1}(B))\cup(A^c\cap Y^{-1}(B)).
\end{split}
\]Since $X$ and $Y$ are random variables, $X^{-1}(B)\in\mathcal{F}$ and $Y^{-1}(B)\in\mathcal{F}$. Moreover, since $\mathcal{F}$ is an algebra, we conclude that\[
(A\cap X^{-1}(B))\cup(A^c\cap Y^{-1}(B))\in\mathcal{F}.
\]
\end{proof}

\section{Exercise 1.2.2: Bounds of the Standard Normal Distribution Function}
Let $\chi$ have the standard normal distribution. Use Theorem~1.2.6 to get upper and lower bounds on $P(\chi\geqslant4)$.
\begin{proof}[Answer]
Letting $x=4$ in Theorem~1.2.6 of \cite{durrett2019probability} gives\[
(\frac{1}{4}-\frac{1}{4^3})e^{-8}\leqslant\int_4^\infty\exp(-\frac{y^2}{2})dy\leqslant\frac{1}{4}e^{-8}.
\]Since\[
P(\chi\geqslant4)=\frac{1}{\sqrt{2\pi}}\int_4^\infty\exp(-\frac{y^2}{2})dy,
\]we conclude that\[
\frac{1}{\sqrt{2\pi}}(\frac{1}{4}-\frac{1}{4^3})e^{-8}\leqslant P(\chi\geqslant4)\leqslant\frac{1}{\sqrt{2\pi}}\cdot\frac{1}{4}\cdot e^{-8}.
\]
\end{proof}

\section{Exercise 1.2.3: A Distribution Function has at most Countable Discontinuities}
Show that a distribution function has at most countably many discontinuities.
\begin{proof}
Pick any distribution function $F$, then, by Theorem~1.2.1 of \cite{durrett2019probability}, $F$ is monotonically increasing on $\mathbb{R}$. Theorem~4.30 of \cite{rudin1976principles} then implies that the set of points of $\mathbb{R}$ at which $F$ is discontinuous is at most countable.
\end{proof}

\section{Exercise 1.3.1: $\sigma(\mathcal{A})=\mathcal{S}$ Implies $\sigma(X^{-1}(\mathcal{A}))=\sigma(X)$}
Show that if $\mathcal{A}$ generates $\mathcal{S}$, then $X^{-1}(\mathcal{A})=\{\{X\in A\}:A\in\mathcal{A}\}$ generates $\sigma(X)=\{\{X\in B\}:B\in\mathcal{S}\}$.
\begin{proof}
Note that,
\begin{gather*}
\sigma(\mathcal{A})=\mathcal{S}\text{ and }\mathcal{A}\subset\sigma(\mathcal{A})\\
\ArrowBetweenLines[\Downarrow]
\mathcal{A}\subset\mathcal{S}\\
\ArrowBetweenLines[\Downarrow]
X^{-1}(\mathcal{A})\subset X^{-1}(\mathcal{S})\\
\ArrowBetweenLines[\Downarrow]
\sigma(X^{-1}(\mathcal{A}))\subset\sigma(X^{-1}(\mathcal{S})).
\end{gather*}
By the definition of $\sigma(X)$, $\sigma(X^{-1}(\mathcal{S}))=\sigma(X)$. Hence we conclude that\[
\sigma(X^{-1}(\mathcal{A}))\subset\sigma(X).
\]

Next, we prove that $\sigma(X)\subset\sigma(X^{-1}(\mathcal{A}))$. Define\[
\mathcal{S}'\equiv\{B\in\mathcal{S}:X^{-1}(B)\in\sigma(X^{-1}(\mathcal{A}))\}.
\]Then,
\begin{gather*}
X^{-1}(\mathcal{A})\subset\sigma(X^{-1}(\mathcal{A}))\\
\ArrowBetweenLines[\Downarrow]
\mathcal{A}\subset\mathcal{S}'\\
\ArrowBetweenLines[\Downarrow]
\sigma(\mathcal{A})\subset\sigma(\mathcal{S}')\\
\ArrowBetweenLines[\Downarrow]
\mathcal{S}\subset\mathcal{S}'.
\end{gather*}
Since $\mathcal{S}'\subset\mathcal{S}$, we conclude that $\mathcal{S}=\mathcal{S}'$. Hence, by the definition of $\sigma(X)$, we conclude that\[
\sigma(X)\subset\sigma(X^{-1}(\mathcal{A})).
\]
\end{proof}

\section{Exercise 1.3.7: The Smallest Class of Simple Functions Closed under Limits}
A function $\phi:\Omega\to\mathbb{R}$ is said to be simple if\[
\phi(\omega)=\sum_{m=1}^nc_m1_{A_m}(\omega)
\]where $c_m$ are real numbers and $A_m\in\mathcal{F}$. Show that the class of $\mathcal{F}$ measurable functions is the smallest class containing the simple functions and closed under pointwise limits.
\begin{proof}
Since $\phi(\omega)$ is $\mathcal{F}$ measurable if and only if the sets $A_1,\dots,A_n$ are $\mathcal{F}$ measurable, it is clear that the class of $\mathcal{F}$ measurable functions contains all simple functions.

Pick any $\mathcal{F}$ measurable function $X$. Then Theorem~11.20 of \cite{rudin1976principles} implies the existence of a sequence $\{s_n\}$ of simple functions such that\[
\lim_{n\to\infty}s_n(\omega)=f(\omega),\quad\forall\omega\in\Omega.
\]Hence $X$ is a member of any class containing the simple functions and closed under pointwise limits. Therefore the class of $\mathcal{F}$ measurable functions is a subset of any class containing the simple functions and closed under pointwise limits.
\end{proof}


\backmatter
\bibliographystyle{amsalpha}
\bibliography{../../bibliography.bib}
\end{document}
